% !TEX root = ../main.tex
% ─────────────────────────────────────────────────────────────────────────────
% Section: Introduction
% Included from main.tex with % !TEX root = ../main.tex
% ─────────────────────────────────────────────────────────────────────────────
% Section: Introduction
% Included from main.tex with % !TEX root = ../main.tex
% ─────────────────────────────────────────────────────────────────────────────
% Section: Introduction
% Included from main.tex with % !TEX root = ../main.tex
% ─────────────────────────────────────────────────────────────────────────────
% Section: Introduction
% Included from main.tex with \input{sections/introduction}
% ─────────────────────────────────────────────────────────────────────────────

\label{sec:introduction}

Location-based services (LBS) have become essential to modern commercial, logistical, and emergency operations. While Global Navigation Satellite Systems (GNSS) deliver reliable sub-metre positioning in open outdoor environments, their microwave signals are heavily attenuated, scattered, and blocked by structural building materials such as reinforced concrete, tinted glass, and steel frames~\cite{torres2014ujiindoorloc}. Consequently, satellite-based localization fails in indoor environments where humans spend upwards of 80\% of their daily lives.

To bridge this operational divide, Indoor Positioning Systems (IPS) leverage opportunistic terrestrial radio-frequency (RF) signals. Among available technologies—including Bluetooth Low Energy (BLE), Ultra-Wideband (UWB), and cellular signals—Wireless Local Area Network (WLAN / Wi-Fi) infrastructure stands out due to its ubiquitous presence across campuses, hospitals, transit terminals, and commercial complexes~\cite{feng2022survey}. Utilizing existing IEEE 802.11 access points (APs) avoids the substantial capital cost of deploying specialized hardware anchor nodes.

Rather than relying on explicit radio propagation modeling (such as log-distance path loss or ray tracing)—which requires knowing exact AP coordinates, antenna patterns, and dynamic indoor floor layouts—\emph{fingerprinting} reframes localization as an empirical pattern-recognition problem~\cite{torres2014ujiindoorloc, torres2015comprehensive}. During an offline survey phase, received signal strength indication (RSSI) values from all detectable APs are cataloged at surveyed reference locations to form a radio map. During online operation, an incoming real-time RSSI vector observed by a mobile device is matched against the database or fed to a predictive model to estimate user position.

Despite its practicality, Wi-Fi RSSI fingerprinting presents three fundamental machine-learning challenges:
\begin{enumerate}
    \item \textbf{High Dimensionality and Extreme Sparsity:} In enterprise environments with hundreds of deployed APs, any single scan detects only a small subset of visible transmitters (typically 3--5\%), leaving over 95\% of feature dimensions unobserved~\cite{torres2014ujiindoorloc}.
    \item \textbf{Complex Propagation Dynamics and Device Heterogeneity:} Multipath fading, temporal signal fluctuation, human shadowing, and antenna gain differences across mobile phone vendors distort recorded decibel levels, introducing substantial noise~\cite{torres2015comprehensive}.
    \item \textbf{Coupled Hierarchical Spatial Targets:} An indoor location is not merely a pair of continuous metric coordinates; it possesses discrete topological context—specifically, which building and floor the user occupies. Misclassifying a floor or building results in severe operational failure (e.g., dispatching emergency responders to the wrong level), even if the 2D Euclidean distance error appears deceptively small.
\end{enumerate}

\subsection{Research Gap and Motivation}
In prevailing literature, indoor positioning targets are predominantly addressed in isolation:
\begin{itemize}
    \item \emph{Decoupled Single-Task Models:} Building and floor classification are frequently trained as distinct classifiers (e.g., using $k$-NN, Support Vector Machines, or Random Forests~\cite{salamah2016enhanced}), while metric coordinate regression is handled by completely independent models. Under this decoupled paradigm, the coordinate regression model cannot leverage the rich inductive biases learned by building and floor classifiers.
    \item \emph{Cascaded Hierarchical Pipelines:} Several architectures (e.g.,~\cite{kim2018scalable}) employ a cascading pipeline: a coarse classifier selects a building, a subsequent classifier selects a floor, and a building/floor-specific sub-network regresses $(x, y)$ coordinates. While physically intuitive, cascading systems suffer from catastrophic \emph{error propagation}: an error in the initial building or floor stage irrevocably ruins downstream coordinate estimation.
    \item \emph{Lack of Controlled End-to-End Multitask Benchmarks:} While deep neural networks have demonstrated strong feature extraction capabilities~\cite{nowicki2017low, song2019cnnloc}, there remains a lack of controlled, unified comparisons on standardized public benchmarks that benchmark classical ML against both single-task and unified multi-task deep architectures on identical input representations.
\end{itemize}

Building upon the theoretical foundations of Multitask Learning (MTL) articulated by Caruana~\cite{caruana1997multitask}, related tasks trained in parallel with a shared internal representation benefit from inductive transfer mechanisms: \emph{data amplification}, \emph{eavesdropping}, \emph{attention focusing}, and \emph{representation bias}. In indoor positioning, learning coarse discrete boundaries (buildings and floors) forces hidden representations to encode macro-scale spatial geometry, which can act as a regularizer for fine-grained coordinate regression.

\subsection{Research Questions}
This investigation is guided by four central research questions:
\begin{itemize}
    \item[\textbf{RQ1:}] \emph{Model Baseline Comparison.} How do classical machine learning baselines ($k$-NN, Random Forest, XGBoost) compare against deep neural network architectures across discrete classification (building, floor) and continuous regression (coordinates)?
    \item[\textbf{RQ2:}] \emph{Inductive Transfer in Multitask Learning.} Does a unified hard-parameter-sharing multitask network outperform isolated single-task models in 2D coordinate regression by jointly optimizing building, floor, and coordinate objectives?
    \item[\textbf{RQ3:}] \emph{Metric Discrepancy.} How do discrete classification accuracy and continuous mean positioning error (MPE) interact, and what trade-offs emerge when balancing loss functions between heterogeneous task heads?
    \item[\textbf{RQ4:}] \emph{Robustness to Signal Sparsity.} How resilient are shared representations to extreme unobserved-signal sparsity compared to instance-based and ensemble baselines?
\end{itemize}

\subsection{Research Objectives and Expected Contributions}
To address these research questions, this project undertakes the following specific objectives:
\begin{enumerate}
    \item \textbf{Comprehensive Literature and Theoretical Review:} Systematically review existing classical ML and deep learning methods for Wi-Fi positioning, establishing mathematical formulations for multi-objective loss optimization.
    \item \textbf{Controlled Benchmark Framework:} Construct a standardized evaluation pipeline on the UJIIndoorLoc public benchmark~\cite{torres2014ujiindoorloc}, ensuring zero validation leakage and rigorous comparison across $k$-NN, Random Forest, XGBoost, Single-Task Deep Networks, and Multitask Deep Networks.
    \item \textbf{Unified Multitask Architecture Design:} Design and implement a hard-parameter-sharing neural network capable of simultaneous multi-output inference (building, floor, longitude, latitude) from raw preprocessed 520-dimensional RSSI fingerprints.
    \item \textbf{Empirical Analysis of Multi-Objective Balancing:} Investigate loss weighting strategies ($\lambda_b, \lambda_f, \lambda_c$) to avoid negative transfer and harmonize categorical cross-entropy with continuous metric regression.
\end{enumerate}

\subsection{Paper Organization}
The remainder of this report is organized as follows: Section~\ref{sec:related-work} reviews related work, comparative taxonomies, and theoretical MTL foundations. Section~\ref{sec:dataset-eda} details the UJIIndoorLoc benchmark and exploratory data analysis. Section~\ref{sec:feature-engineering} outlines the preprocessing protocol and missing-signal handling. Section~\ref{sec:methodology} describes the model architectures, loss formulations, and training strategies. Section~\ref{sec:results} presents comparative empirical results, ablation studies, and discussion. Finally, Section~\ref{sec:conclusion} concludes the paper with future research directions.

% ─────────────────────────────────────────────────────────────────────────────

\label{sec:introduction}

Location-based services (LBS) have become essential to modern commercial, logistical, and emergency operations. While Global Navigation Satellite Systems (GNSS) deliver reliable sub-metre positioning in open outdoor environments, their microwave signals are heavily attenuated, scattered, and blocked by structural building materials such as reinforced concrete, tinted glass, and steel frames~\cite{torres2014ujiindoorloc}. Consequently, satellite-based localization fails in indoor environments where humans spend upwards of 80\% of their daily lives.

To bridge this operational divide, Indoor Positioning Systems (IPS) leverage opportunistic terrestrial radio-frequency (RF) signals. Among available technologies—including Bluetooth Low Energy (BLE), Ultra-Wideband (UWB), and cellular signals—Wireless Local Area Network (WLAN / Wi-Fi) infrastructure stands out due to its ubiquitous presence across campuses, hospitals, transit terminals, and commercial complexes~\cite{feng2022survey}. Utilizing existing IEEE 802.11 access points (APs) avoids the substantial capital cost of deploying specialized hardware anchor nodes.

Rather than relying on explicit radio propagation modeling (such as log-distance path loss or ray tracing)—which requires knowing exact AP coordinates, antenna patterns, and dynamic indoor floor layouts—\emph{fingerprinting} reframes localization as an empirical pattern-recognition problem~\cite{torres2014ujiindoorloc, torres2015comprehensive}. During an offline survey phase, received signal strength indication (RSSI) values from all detectable APs are cataloged at surveyed reference locations to form a radio map. During online operation, an incoming real-time RSSI vector observed by a mobile device is matched against the database or fed to a predictive model to estimate user position.

Despite its practicality, Wi-Fi RSSI fingerprinting presents three fundamental machine-learning challenges:
\begin{enumerate}
    \item \textbf{High Dimensionality and Extreme Sparsity:} In enterprise environments with hundreds of deployed APs, any single scan detects only a small subset of visible transmitters (typically 3--5\%), leaving over 95\% of feature dimensions unobserved~\cite{torres2014ujiindoorloc}.
    \item \textbf{Complex Propagation Dynamics and Device Heterogeneity:} Multipath fading, temporal signal fluctuation, human shadowing, and antenna gain differences across mobile phone vendors distort recorded decibel levels, introducing substantial noise~\cite{torres2015comprehensive}.
    \item \textbf{Coupled Hierarchical Spatial Targets:} An indoor location is not merely a pair of continuous metric coordinates; it possesses discrete topological context—specifically, which building and floor the user occupies. Misclassifying a floor or building results in severe operational failure (e.g., dispatching emergency responders to the wrong level), even if the 2D Euclidean distance error appears deceptively small.
\end{enumerate}

\subsection{Research Gap and Motivation}
In prevailing literature, indoor positioning targets are predominantly addressed in isolation:
\begin{itemize}
    \item \emph{Decoupled Single-Task Models:} Building and floor classification are frequently trained as distinct classifiers (e.g., using $k$-NN, Support Vector Machines, or Random Forests~\cite{salamah2016enhanced}), while metric coordinate regression is handled by completely independent models. Under this decoupled paradigm, the coordinate regression model cannot leverage the rich inductive biases learned by building and floor classifiers.
    \item \emph{Cascaded Hierarchical Pipelines:} Several architectures (e.g.,~\cite{kim2018scalable}) employ a cascading pipeline: a coarse classifier selects a building, a subsequent classifier selects a floor, and a building/floor-specific sub-network regresses $(x, y)$ coordinates. While physically intuitive, cascading systems suffer from catastrophic \emph{error propagation}: an error in the initial building or floor stage irrevocably ruins downstream coordinate estimation.
    \item \emph{Lack of Controlled End-to-End Multitask Benchmarks:} While deep neural networks have demonstrated strong feature extraction capabilities~\cite{nowicki2017low, song2019cnnloc}, there remains a lack of controlled, unified comparisons on standardized public benchmarks that benchmark classical ML against both single-task and unified multi-task deep architectures on identical input representations.
\end{itemize}

Building upon the theoretical foundations of Multitask Learning (MTL) articulated by Caruana~\cite{caruana1997multitask}, related tasks trained in parallel with a shared internal representation benefit from inductive transfer mechanisms: \emph{data amplification}, \emph{eavesdropping}, \emph{attention focusing}, and \emph{representation bias}. In indoor positioning, learning coarse discrete boundaries (buildings and floors) forces hidden representations to encode macro-scale spatial geometry, which can act as a regularizer for fine-grained coordinate regression.

\subsection{Research Questions}
This investigation is guided by four central research questions:
\begin{itemize}
    \item[\textbf{RQ1:}] \emph{Model Baseline Comparison.} How do classical machine learning baselines ($k$-NN, Random Forest, XGBoost) compare against deep neural network architectures across discrete classification (building, floor) and continuous regression (coordinates)?
    \item[\textbf{RQ2:}] \emph{Inductive Transfer in Multitask Learning.} Does a unified hard-parameter-sharing multitask network outperform isolated single-task models in 2D coordinate regression by jointly optimizing building, floor, and coordinate objectives?
    \item[\textbf{RQ3:}] \emph{Metric Discrepancy.} How do discrete classification accuracy and continuous mean positioning error (MPE) interact, and what trade-offs emerge when balancing loss functions between heterogeneous task heads?
    \item[\textbf{RQ4:}] \emph{Robustness to Signal Sparsity.} How resilient are shared representations to extreme unobserved-signal sparsity compared to instance-based and ensemble baselines?
\end{itemize}

\subsection{Research Objectives and Expected Contributions}
To address these research questions, this project undertakes the following specific objectives:
\begin{enumerate}
    \item \textbf{Comprehensive Literature and Theoretical Review:} Systematically review existing classical ML and deep learning methods for Wi-Fi positioning, establishing mathematical formulations for multi-objective loss optimization.
    \item \textbf{Controlled Benchmark Framework:} Construct a standardized evaluation pipeline on the UJIIndoorLoc public benchmark~\cite{torres2014ujiindoorloc}, ensuring zero validation leakage and rigorous comparison across $k$-NN, Random Forest, XGBoost, Single-Task Deep Networks, and Multitask Deep Networks.
    \item \textbf{Unified Multitask Architecture Design:} Design and implement a hard-parameter-sharing neural network capable of simultaneous multi-output inference (building, floor, longitude, latitude) from raw preprocessed 520-dimensional RSSI fingerprints.
    \item \textbf{Empirical Analysis of Multi-Objective Balancing:} Investigate loss weighting strategies ($\lambda_b, \lambda_f, \lambda_c$) to avoid negative transfer and harmonize categorical cross-entropy with continuous metric regression.
\end{enumerate}

\subsection{Paper Organization}
The remainder of this report is organized as follows: Section~\ref{sec:related-work} reviews related work, comparative taxonomies, and theoretical MTL foundations. Section~\ref{sec:dataset-eda} details the UJIIndoorLoc benchmark and exploratory data analysis. Section~\ref{sec:feature-engineering} outlines the preprocessing protocol and missing-signal handling. Section~\ref{sec:methodology} describes the model architectures, loss formulations, and training strategies. Section~\ref{sec:results} presents comparative empirical results, ablation studies, and discussion. Finally, Section~\ref{sec:conclusion} concludes the paper with future research directions.

% ─────────────────────────────────────────────────────────────────────────────

\label{sec:introduction}

Location-based services (LBS) have become essential to modern commercial, logistical, and emergency operations. While Global Navigation Satellite Systems (GNSS) deliver reliable sub-metre positioning in open outdoor environments, their microwave signals are heavily attenuated, scattered, and blocked by structural building materials such as reinforced concrete, tinted glass, and steel frames~\cite{torres2014ujiindoorloc}. Consequently, satellite-based localization fails in indoor environments where humans spend upwards of 80\% of their daily lives.

To bridge this operational divide, Indoor Positioning Systems (IPS) leverage opportunistic terrestrial radio-frequency (RF) signals. Among available technologies—including Bluetooth Low Energy (BLE), Ultra-Wideband (UWB), and cellular signals—Wireless Local Area Network (WLAN / Wi-Fi) infrastructure stands out due to its ubiquitous presence across campuses, hospitals, transit terminals, and commercial complexes~\cite{feng2022survey}. Utilizing existing IEEE 802.11 access points (APs) avoids the substantial capital cost of deploying specialized hardware anchor nodes.

Rather than relying on explicit radio propagation modeling (such as log-distance path loss or ray tracing)—which requires knowing exact AP coordinates, antenna patterns, and dynamic indoor floor layouts—\emph{fingerprinting} reframes localization as an empirical pattern-recognition problem~\cite{torres2014ujiindoorloc, torres2015comprehensive}. During an offline survey phase, received signal strength indication (RSSI) values from all detectable APs are cataloged at surveyed reference locations to form a radio map. During online operation, an incoming real-time RSSI vector observed by a mobile device is matched against the database or fed to a predictive model to estimate user position.

Despite its practicality, Wi-Fi RSSI fingerprinting presents three fundamental machine-learning challenges:
\begin{enumerate}
    \item \textbf{High Dimensionality and Extreme Sparsity:} In enterprise environments with hundreds of deployed APs, any single scan detects only a small subset of visible transmitters (typically 3--5\%), leaving over 95\% of feature dimensions unobserved~\cite{torres2014ujiindoorloc}.
    \item \textbf{Complex Propagation Dynamics and Device Heterogeneity:} Multipath fading, temporal signal fluctuation, human shadowing, and antenna gain differences across mobile phone vendors distort recorded decibel levels, introducing substantial noise~\cite{torres2015comprehensive}.
    \item \textbf{Coupled Hierarchical Spatial Targets:} An indoor location is not merely a pair of continuous metric coordinates; it possesses discrete topological context—specifically, which building and floor the user occupies. Misclassifying a floor or building results in severe operational failure (e.g., dispatching emergency responders to the wrong level), even if the 2D Euclidean distance error appears deceptively small.
\end{enumerate}

\subsection{Research Gap and Motivation}
In prevailing literature, indoor positioning targets are predominantly addressed in isolation:
\begin{itemize}
    \item \emph{Decoupled Single-Task Models:} Building and floor classification are frequently trained as distinct classifiers (e.g., using $k$-NN, Support Vector Machines, or Random Forests~\cite{salamah2016enhanced}), while metric coordinate regression is handled by completely independent models. Under this decoupled paradigm, the coordinate regression model cannot leverage the rich inductive biases learned by building and floor classifiers.
    \item \emph{Cascaded Hierarchical Pipelines:} Several architectures (e.g.,~\cite{kim2018scalable}) employ a cascading pipeline: a coarse classifier selects a building, a subsequent classifier selects a floor, and a building/floor-specific sub-network regresses $(x, y)$ coordinates. While physically intuitive, cascading systems suffer from catastrophic \emph{error propagation}: an error in the initial building or floor stage irrevocably ruins downstream coordinate estimation.
    \item \emph{Lack of Controlled End-to-End Multitask Benchmarks:} While deep neural networks have demonstrated strong feature extraction capabilities~\cite{nowicki2017low, song2019cnnloc}, there remains a lack of controlled, unified comparisons on standardized public benchmarks that benchmark classical ML against both single-task and unified multi-task deep architectures on identical input representations.
\end{itemize}

Building upon the theoretical foundations of Multitask Learning (MTL) articulated by Caruana~\cite{caruana1997multitask}, related tasks trained in parallel with a shared internal representation benefit from inductive transfer mechanisms: \emph{data amplification}, \emph{eavesdropping}, \emph{attention focusing}, and \emph{representation bias}. In indoor positioning, learning coarse discrete boundaries (buildings and floors) forces hidden representations to encode macro-scale spatial geometry, which can act as a regularizer for fine-grained coordinate regression.

\subsection{Research Questions}
This investigation is guided by four central research questions:
\begin{itemize}
    \item[\textbf{RQ1:}] \emph{Model Baseline Comparison.} How do classical machine learning baselines ($k$-NN, Random Forest, XGBoost) compare against deep neural network architectures across discrete classification (building, floor) and continuous regression (coordinates)?
    \item[\textbf{RQ2:}] \emph{Inductive Transfer in Multitask Learning.} Does a unified hard-parameter-sharing multitask network outperform isolated single-task models in 2D coordinate regression by jointly optimizing building, floor, and coordinate objectives?
    \item[\textbf{RQ3:}] \emph{Metric Discrepancy.} How do discrete classification accuracy and continuous mean positioning error (MPE) interact, and what trade-offs emerge when balancing loss functions between heterogeneous task heads?
    \item[\textbf{RQ4:}] \emph{Robustness to Signal Sparsity.} How resilient are shared representations to extreme unobserved-signal sparsity compared to instance-based and ensemble baselines?
\end{itemize}

\subsection{Research Objectives and Expected Contributions}
To address these research questions, this project undertakes the following specific objectives:
\begin{enumerate}
    \item \textbf{Comprehensive Literature and Theoretical Review:} Systematically review existing classical ML and deep learning methods for Wi-Fi positioning, establishing mathematical formulations for multi-objective loss optimization.
    \item \textbf{Controlled Benchmark Framework:} Construct a standardized evaluation pipeline on the UJIIndoorLoc public benchmark~\cite{torres2014ujiindoorloc}, ensuring zero validation leakage and rigorous comparison across $k$-NN, Random Forest, XGBoost, Single-Task Deep Networks, and Multitask Deep Networks.
    \item \textbf{Unified Multitask Architecture Design:} Design and implement a hard-parameter-sharing neural network capable of simultaneous multi-output inference (building, floor, longitude, latitude) from raw preprocessed 520-dimensional RSSI fingerprints.
    \item \textbf{Empirical Analysis of Multi-Objective Balancing:} Investigate loss weighting strategies ($\lambda_b, \lambda_f, \lambda_c$) to avoid negative transfer and harmonize categorical cross-entropy with continuous metric regression.
\end{enumerate}

\subsection{Paper Organization}
The remainder of this report is organized as follows: Section~\ref{sec:related-work} reviews related work, comparative taxonomies, and theoretical MTL foundations. Section~\ref{sec:dataset-eda} details the UJIIndoorLoc benchmark and exploratory data analysis. Section~\ref{sec:feature-engineering} outlines the preprocessing protocol and missing-signal handling. Section~\ref{sec:methodology} describes the model architectures, loss formulations, and training strategies. Section~\ref{sec:results} presents comparative empirical results, ablation studies, and discussion. Finally, Section~\ref{sec:conclusion} concludes the paper with future research directions.

% ─────────────────────────────────────────────────────────────────────────────

\label{sec:introduction}

Location-based services (LBS) have become essential to modern commercial, logistical, and emergency operations. While Global Navigation Satellite Systems (GNSS) deliver reliable sub-metre positioning in open outdoor environments, their microwave signals are heavily attenuated, scattered, and blocked by structural building materials such as reinforced concrete, tinted glass, and steel frames~\cite{torres2014ujiindoorloc}. Consequently, satellite-based localization fails in indoor environments where humans spend upwards of 80\% of their daily lives.

To bridge this operational divide, Indoor Positioning Systems (IPS) leverage opportunistic terrestrial radio-frequency (RF) signals. Among available technologies—including Bluetooth Low Energy (BLE), Ultra-Wideband (UWB), and cellular signals—Wireless Local Area Network (WLAN / Wi-Fi) infrastructure stands out due to its ubiquitous presence across campuses, hospitals, transit terminals, and commercial complexes~\cite{feng2022survey}. Utilizing existing IEEE 802.11 access points (APs) avoids the substantial capital cost of deploying specialized hardware anchor nodes.

Rather than relying on explicit radio propagation modeling (such as log-distance path loss or ray tracing)—which requires knowing exact AP coordinates, antenna patterns, and dynamic indoor floor layouts—\emph{fingerprinting} reframes localization as an empirical pattern-recognition problem~\cite{torres2014ujiindoorloc, torres2015comprehensive}. During an offline survey phase, received signal strength indication (RSSI) values from all detectable APs are cataloged at surveyed reference locations to form a radio map. During online operation, an incoming real-time RSSI vector observed by a mobile device is matched against the database or fed to a predictive model to estimate user position.

Despite its practicality, Wi-Fi RSSI fingerprinting presents three fundamental machine-learning challenges:
\begin{enumerate}
    \item \textbf{High Dimensionality and Extreme Sparsity:} In enterprise environments with hundreds of deployed APs, any single scan detects only a small subset of visible transmitters (typically 3--5\%), leaving over 95\% of feature dimensions unobserved~\cite{torres2014ujiindoorloc}.
    \item \textbf{Complex Propagation Dynamics and Device Heterogeneity:} Multipath fading, temporal signal fluctuation, human shadowing, and antenna gain differences across mobile phone vendors distort recorded decibel levels, introducing substantial noise~\cite{torres2015comprehensive}.
    \item \textbf{Coupled Hierarchical Spatial Targets:} An indoor location is not merely a pair of continuous metric coordinates; it possesses discrete topological context—specifically, which building and floor the user occupies. Misclassifying a floor or building results in severe operational failure (e.g., dispatching emergency responders to the wrong level), even if the 2D Euclidean distance error appears deceptively small.
\end{enumerate}

\subsection{Research Gap and Motivation}
In prevailing literature, indoor positioning targets are predominantly addressed in isolation:
\begin{itemize}
    \item \emph{Decoupled Single-Task Models:} Building and floor classification are frequently trained as distinct classifiers (e.g., using $k$-NN, Support Vector Machines, or Random Forests~\cite{salamah2016enhanced}), while metric coordinate regression is handled by completely independent models. Under this decoupled paradigm, the coordinate regression model cannot leverage the rich inductive biases learned by building and floor classifiers.
    \item \emph{Cascaded Hierarchical Pipelines:} Several architectures (e.g.,~\cite{kim2018scalable}) employ a cascading pipeline: a coarse classifier selects a building, a subsequent classifier selects a floor, and a building/floor-specific sub-network regresses $(x, y)$ coordinates. While physically intuitive, cascading systems suffer from catastrophic \emph{error propagation}: an error in the initial building or floor stage irrevocably ruins downstream coordinate estimation.
    \item \emph{Lack of Controlled End-to-End Multitask Benchmarks:} While deep neural networks have demonstrated strong feature extraction capabilities~\cite{nowicki2017low, song2019cnnloc}, there remains a lack of controlled, unified comparisons on standardized public benchmarks that benchmark classical ML against both single-task and unified multi-task deep architectures on identical input representations.
\end{itemize}

Building upon the theoretical foundations of Multitask Learning (MTL) articulated by Caruana~\cite{caruana1997multitask}, related tasks trained in parallel with a shared internal representation benefit from inductive transfer mechanisms: \emph{data amplification}, \emph{eavesdropping}, \emph{attention focusing}, and \emph{representation bias}. In indoor positioning, learning coarse discrete boundaries (buildings and floors) forces hidden representations to encode macro-scale spatial geometry, which can act as a regularizer for fine-grained coordinate regression.

\subsection{Research Questions}
This investigation is guided by four central research questions:
\begin{itemize}
    \item[\textbf{RQ1:}] \emph{Model Baseline Comparison.} How do classical machine learning baselines ($k$-NN, Random Forest, XGBoost) compare against deep neural network architectures across discrete classification (building, floor) and continuous regression (coordinates)?
    \item[\textbf{RQ2:}] \emph{Inductive Transfer in Multitask Learning.} Does a unified hard-parameter-sharing multitask network outperform isolated single-task models in 2D coordinate regression by jointly optimizing building, floor, and coordinate objectives?
    \item[\textbf{RQ3:}] \emph{Metric Discrepancy.} How do discrete classification accuracy and continuous mean positioning error (MPE) interact, and what trade-offs emerge when balancing loss functions between heterogeneous task heads?
    \item[\textbf{RQ4:}] \emph{Robustness to Signal Sparsity.} How resilient are shared representations to extreme unobserved-signal sparsity compared to instance-based and ensemble baselines?
\end{itemize}

\subsection{Research Objectives and Expected Contributions}
To address these research questions, this project undertakes the following specific objectives:
\begin{enumerate}
    \item \textbf{Comprehensive Literature and Theoretical Review:} Systematically review existing classical ML and deep learning methods for Wi-Fi positioning, establishing mathematical formulations for multi-objective loss optimization.
    \item \textbf{Controlled Benchmark Framework:} Construct a standardized evaluation pipeline on the UJIIndoorLoc public benchmark~\cite{torres2014ujiindoorloc}, ensuring zero validation leakage and rigorous comparison across $k$-NN, Random Forest, XGBoost, Single-Task Deep Networks, and Multitask Deep Networks.
    \item \textbf{Unified Multitask Architecture Design:} Design and implement a hard-parameter-sharing neural network capable of simultaneous multi-output inference (building, floor, longitude, latitude) from raw preprocessed 520-dimensional RSSI fingerprints.
    \item \textbf{Empirical Analysis of Multi-Objective Balancing:} Investigate loss weighting strategies ($\lambda_b, \lambda_f, \lambda_c$) to avoid negative transfer and harmonize categorical cross-entropy with continuous metric regression.
\end{enumerate}

\subsection{Paper Organization}
The remainder of this report is organized as follows: Section~\ref{sec:related-work} reviews related work, comparative taxonomies, and theoretical MTL foundations. Section~\ref{sec:dataset-eda} details the UJIIndoorLoc benchmark and exploratory data analysis. Section~\ref{sec:feature-engineering} outlines the preprocessing protocol and missing-signal handling. Section~\ref{sec:methodology} describes the model architectures, loss formulations, and training strategies. Section~\ref{sec:results} presents comparative empirical results, ablation studies, and discussion. Finally, Section~\ref{sec:conclusion} concludes the paper with future research directions.
